\section{String Transform}
\label{sec:cses-string-transform}

\problemstatement{The Burrows--Wheeler transform (BWT) of a string sorts
all rotations of the string (with an appended sentinel \texttt{\#} that is
lexicographically smallest) and outputs the last column. Given a BWT
string, reconstruct the original string.}

\begin{patternbox}
\emph{``Invert the Burrows--Wheeler transform''} is done with the
\textbf{LF-mapping}: from the last column alone you can recover the first
column (it is just the last column sorted), and a rank-based mapping links
each character to its predecessor, letting you walk the original string out
one character at a time.
\end{patternbox}

\subsection*{LF-mapping in one idea}

Let $L$ be the given last column. The first column $F$ is $L$ sorted. The
\textbf{LF-mapping} exploits a key invariant: the $k$-th occurrence of a
character in $L$ corresponds to the $k$-th occurrence of that same
character in $F$. Concretely, for a row $i$, the character $L[i]$ maps to
row $\textit{start}[L[i]]+\textit{rank}_i$, where $\textit{start}[c]$ is
the first row in $F$ beginning with $c$ and $\textit{rank}_i$ is how many
$L[i]$ appeared before index $i$. Following this mapping from the sentinel
row spells the original string in reverse.

\begin{ideabox}[Same Rank in Last and First Columns]
Because equal characters keep their relative order between the sorted-first
and last columns, the LF-mapping is exact: it reconstructs which character
precedes each one in the original. Iterating it $n$ times recovers the
whole string with no ambiguity.
\end{ideabox}

\subsection*{Algorithm}

\begin{enumerate}
  \item Count characters of $L$; compute $\textit{start}[c]$ = number of
        characters smaller than $c$.
  \item Compute $\textit{rank}_i$ for each row and the LF map
        $\textit{lf}[i]=\textit{start}[L[i]]+\textit{rank}_i$.
  \item Start at the sentinel row (row $0$ of $F$), emit $L[\text{row}]$,
        follow $\textit{lf}$; after $n$ steps reverse to get the original.
\end{enumerate}

\subsection*{Dry run}

\dryrun{Conceptually, $L$'s sorted copy is $F$; equal letters keep order.}

\begin{itemize}
  \item From the sentinel row, each LF step jumps to the row whose first
        character is the current $L$ character, at the matching rank.
  \item Emitting characters along this chain and reversing reproduces the
        original string ending in \texttt{\#}.
\end{itemize}

\subsection*{C++ implementation}

\begin{lstlisting}
#include <bits/stdc++.h>
using namespace std;

int main() {
    string L;
    cin >> L;                                    // BWT (last column), contains '#'
    int n = L.size();

    vector<int> cnt(128, 0);
    for (char c : L) cnt[c]++;
    vector<int> start(128, 0);
    for (int c = 1; c < 128; c++) start[c] = start[c - 1] + cnt[c - 1];

    vector<int> lf(n);
    vector<int> occ(128, 0);
    for (int i = 0; i < n; i++) {
        lf[i] = start[(int)L[i]] + occ[(int)L[i]];
        occ[(int)L[i]]++;
    }

    // sentinel '#' is smallest -> first column row 0 begins with '#'
    string res;
    int row = 0;
    for (int i = 0; i < n; i++) {
        res += L[row];
        row = lf[row];
    }
    reverse(res.begin(), res.end());
    cout << res << "\n";                          // original (ending in '#')
    return 0;
}
\end{lstlisting}

\begin{complexitybox}
\textbf{Time Complexity.} $O(n+\Sigma)$ -- linear counting sort plus a
linear LF walk. \textbf{Space Complexity.} $O(n+\Sigma)$.
\end{complexitybox}

\begin{takeawaybox}
Inverting the BWT hinges on the LF-mapping's rank-preservation between the
first and last columns, letting you unwind the original one character at a
time in linear time. The same rank correspondence underlies FM-index
search used in bioinformatics.
\end{takeawaybox}
